\documentclass[../../../include/open-logic-section]{subfiles}

\begin{document}

\olfileid{sth}{replacement}{limofsize}
\olsection{আকার-সীমার নীতি}

প্রতিস্থাপনের ``অন্তর্নিহিত'' ন্যায্যতা দেওয়ার
সম্ভবত সবচেয়ে প্রচলিত চেষ্টা নিচের ধারণা
থেকে আসে:
\begin{enumerate}
	\item[] \limofsize. যেকোনো বস্তুসমষ্টি একটি
	সেট গঠন করে, যদি বস্তুগুলির সংখ্যা অতিরিক্ত
	বেশি না হয়।
\end{enumerate}
এই নীতি সরাসরি প্রতিস্থাপনকে সমর্থন করে।
কারণ প্রতিস্থাপন দিয়ে গঠিত কোনো সেট
যে সেট থেকে গঠিত, তার চেয়ে বড় হতে পারে না।
নির্দিষ্টভাবে: প্রতিস্থাপন ব্যবহার করে
$\funimage{\tau}{A} = \Setabs{\tau(x)}{x \in A}$
গঠন করলে $\cardle{\funimage{\tau}{A}}{A}$;
তাই $A$-র !!{element}s যদি সেট গঠনের
পক্ষে অতিরিক্ত সংখ্যক না হয়, তবে তাদের
প্রতিচ্ছবিগুলিও $\funimage{\tau}{A}$ গঠনের
পক্ষে অতিরিক্ত সংখ্যক নয়।
% BN-SRC-792 TARGET-CORRECTION-BEGIN
\emph{উৎস-সংশোধন-নোট।} বইয়ের সদস্যসংখ্যার এই তুলনা
প্রতিচ্ছবি থেকে $A$-তে একৈক চিত্রণের অস্তিত্ব বোঝায়। কেবল
$A$ থেকে প্রতিচ্ছবিতে সমাপতিত চিত্রণ থাকায় সেটি নির্বিশেষে
প্রমাণিত হয় না। $A$-তে সুক্রম থাকলে প্রত্যেক পূর্বপ্রতিচ্ছবিতের
ক্ষুদ্রতম সদস্য বেছে একৈক বিপরীত পাই; নির্বাচনের স্বতঃসিদ্ধও
যথেষ্ট। এই বাড়তি অনুমান ছাড়া এখানে নিশ্চিত ফল শুধু
সমাপতন। আকার-সীমার অনানুষ্ঠানিক প্রেরণাকে ওই সমাপতনভিত্তিক
অর্থে পড়তে হবে, অথবা তুলনার জন্য বাড়তি নীতি স্পষ্ট নিতে হবে।
% BN-SRC-792 TARGET-CORRECTION-END

প্রতিস্থাপনের ন্যায্যতায় \limofsize{}
আহ্বানের স্পষ্ট অসুবিধা হলো, \limofsize-র
মতো কোনো নীতি আমরা \emph{এখনও}
স্থাপন করিনি। উপরন্তু,
\crefrange{sth:story::chap}{sth:z::chap}-এ
সেটের সঞ্চয়ী-পর্যায়ক্রমিক ধারণার আখ্যান
বলতে গিয়ে \limofsize-র দিকে কোনো
\emph{ইঙ্গিতও} দিইনি। ঠিক এই কারণেই
বুলস এক সময় লিখেছিলেন: ``সম্ভবত
সিদ্ধান্ত করা যায়, সেটতত্ত্বের ‘পেছনে’
অন্তত দুটি চিন্তা আছে''
\citeyearpar[p.~19]{Boolos1989}।
এক দিকে সেটের সঞ্চয়ী-পর্যায়ক্রমিক
ধারণা $\Z$-কে সমর্থন করার কথা;
অন্য দিকে \limofsize{} প্রতিস্থাপনকে
সমর্থন করার কথা।

তবে সমস্যা শুধু এই নয় যে \limofsize{}
নিয়ে আমরা এ পর্যন্ত \emph{নীরব} ছিলাম।
সমস্যা আরও গভীর: এইভাবে বলা \limofsize{}
সেটের সঞ্চয়ী-পর্যায়ক্রমিক ধারণার সঙ্গে
খাপ খায় না বলেই মনে হয়। কারণ সেট
\emph{পর্যায়ে} গঠিত হয়—এই ধারণার
কোনো উল্লেখ এতে নেই।

এই ``দুই চিন্তা''—সেটের সঞ্চয়ী-পর্যায়ক্রমিক
ধারণা ও \limofsize-র ইতিহাস বিবেচনায়
এতে খুব অবাক হওয়ার নেই। শেষ পর্যন্ত
সেটতাত্ত্বিক কূটাভাসগুলি ঠেকানোর দুটি
ভিন্ন প্রকল্প এদের পেছনে আছে। সঞ্চয়ী-পর্যায়ক্রমিক
ধারণা মোটামুটি বলে রাসেলের কূটাভাস ঠেকে:
\emph{রাসেল-সেট থাকবে বলে আশা করাই উচিত
ছিল না, কারণ সেটি কোনো পর্যায়েই ``গঠিত''
হতো না}। বিপরীতে, \limofsize{} মোটামুটি
বলে রাসেল-সেট বাদ পড়ে: \emph{রাসেল-সেট
থাকবে বলে আশা করাই উচিত ছিল না,
কারণ সেটি অতিরিক্ত বড় হতো}।

এভাবে বললে স্পষ্ট কথাটি বলতেই হয়:
কূটাভাসের জবাব হিসেবে \limofsize-র
\emph{অনেক} বেশি ন্যায্যতা দরকার।
যেমন রাসেলের কূটাভাসের এই রূপটি ভাবুন:
\emph{এমন কোনো পাগ কুকুর নেই, যে ঠিক সেই পাগগুলিকেই
শোঁকে যারা নিজেদের শোঁকে না}
(দেখুন \olref[sth][story][rus]{sec})।
``এমন পাগ নেই কেন?'' জিজ্ঞেস করলে
``সেই পাগকে অত্যধিক পাগ শুঁকতে হতো''
ভালো উত্তর নয়। তবে রাসেল-সেটের
অনস্তিত্বের স্বজ্ঞাগত ব্যাখ্যা হিসেবে
``সেটি অস্তিত্বের পক্ষে অতিরিক্ত বড়
হতো'' কেন ভালো উত্তর হবে?

তাই \limofsize-র ``স্বজ্ঞাগত'' প্রেরণা
আপনার কাছে কিছুটা দুর্বোধ্য মনে হলে
তা অস্বাভাবিক নয়।

\end{document}
